% Generated draft: requires mathematical and source review.
\documentclass[11pt]{article}
\usepackage[margin=25mm]{geometry}
\usepackage{fontspec}
\setmainfont{Latin Modern Roman}
\usepackage{amsmath,amssymb,mathtools,unicode-math}
\setmathfont{Latin Modern Math}
\usepackage{enumitem,needspace,xurl,hyperref}
\setlength{\parindent}{0pt}
\setlength{\parskip}{5pt}
\setlength{\emergencystretch}{3em}
\newcommand{\N}{\mathbb{N}}
\newcommand{\Z}{\mathbb{Z}}
\newcommand{\Q}{\mathbb{Q}}
\newcommand{\R}{\mathbb{R}}
\newcommand{\C}{\mathbb{C}}
\newcommand{\F}{\mathbb{F}}
\newcommand{\A}{\mathbb{A}}
\newcommand{\PP}{\mathbb{P}}
\newcommand{\E}{\mathbb{E}}
\newcommand{\eps}{\varepsilon}
\newcommand{\dd}{\,\mathrm{d}}
\begin{document}
\section*{UnsolvedMath ID: 30004699}
\textbf{Source problem code:} catalogue entry
\subsection{Definitions and mathematical statement}
Here \emph{artha} denotes the meanings and purposes that commentaries assign to mathematical treatises, following the supplied wording; no more specific technical definition is supplied.

\textbf{Research question.} What meanings and purposes, or \emph{artha}, do commentaries give to mathematical treatises, and what mathematical work is required to establish those meanings and purposes?

This is an interpretive research question concerning mathematical texts and their commentaries. The record specifies neither a corpus of treatises and commentaries nor historical, linguistic, or geographical boundaries. It also leaves unspecified what constitutes establishing an interpretation and what kinds of mathematical work are intended. Consequently, the supplied evidence does not determine a quantified mathematical proposition or a precise criterion for a complete answer. These qualifications remain unresolved.
\subsection{Short English statement}
Investigate the meanings and purposes that commentaries assign to mathematical treatises and the mathematical work needed to establish them. The intended corpus and standards of justification require clarification. The supplied curation treats this as an interpretive research question, not a claim with a settled yes-or-no resolution.
\subsection{Sources}
\begin{itemize}
\item Local UnsolvedMath registry, supplied complete source record: permanent ID \(30004699\), title \emph{Meanings and Purposes in Mathematical Commentaries}, collection \emph{Oberwolfach Reports Open Problems}, exact source problem code \texttt{catalogue entry}. No individual proposer is identified; the supplied proposed year is 2021.
\item Supplied report citation: \emph{Mathematics and its Ancient Classics Worldwide: Translations, Appropriations, Reconstructions, Roles} (2022), \emph{Oberwolfach Reports}, pp.\ 1347--1406, DOI 10.4171/owr/2021/26. Supplied source URL and literature reference labelled ``Oberwolfach report'': \url{https://doi.org/10.4171/owr/2021/26}. The background describes Oberwolfach Reports as documenting mathematical workshops and their open-problem discussions.
\item Literature assessment dated \texttt{2026-08-22}: the supplied curation labels the item ``open'' and describes the comparison of meanings and purposes assigned by commentaries as an interpretive research question without a settled yes-or-no resolution. This is supplied dated curation derived from the OWR record, not an independently verified literature assessment. The background requests independent verification of time-sensitive status and cited claims.
\item The supplied registry marks the statement \texttt{corrected\_verified} and reports a manual semantic, typographic, and LaTeX review confirming self-containment and preservation of the target without unsupported additions. That reported review is source metadata; it does not resolve the formulation issues identified here or constitute independent verification in this draft.
\end{itemize}
\end{document}
